% Reviewed mathematical development fragment; not a replacement manuscript.
% Source status is recorded in development-brief.json and the per-direction
% records. Lean candidates are not accepted merely because this prose is proved.

\subsection{Cyclic content and the boundary of the density inequality}

The canonical description separates two sources of density behaviour: the
cyclic core of an atom and the locations at which atoms are attached.  The
following low-rank observation complements the high-rank negative example.
It is a deduction from graph recognition and known theta inequalities,
not a new proof of those inequalities.

\begin{proposition}[The low-rank single-atom boundary]
\label{density-low-rank-atom-boundary}
Let $r\ge3$, and let $A$ be a single canonical obligatory $r$-uniform atom.
If the cycle rank of its connected Levi graph is at most two, then $A$ is
Sidorenko.  Consequently a non-Sidorenko single canonical atom has Levi
cycle rank at least three.
\end{proposition}
\begin{proof}
A single-edge atom has density $p$ and cycle rank zero.  Otherwise the
canonical description gives $A=J^{[r]}$, where $J$ is a finite simple
$2$-connected bipartite graph.  If $n=|V(J)|$ and $m=|E(J)|$, then
\[
 \beta(A)=rm-\bigl(n+(r-1)m\bigr)+1=m-n+1=\beta(J).
\]
Every vertex of $J$ has degree at least two, so $\beta(J)\ge1$.
At rank one, the degree sum forces all degrees to be two; thus $J$ is an
even cycle.  At rank two, the degree excess is
$\sum_v(\deg(v)-2)=2$.  The possibility of a single vertex of degree four
is excluded by $2$-connectivity: deleting it would leave two path
components.  Hence there are exactly two vertices of degree three, and
the other vertices lie on three internally disjoint paths between them.
Bipartiteness forces the three path lengths to have the same parity.

For a symmetric $[0,1]$-valued $r$-kernel $U$ on a finite probability
space, let $M_U$ be its two-coordinate marginal.  Integrating each
edge's distinct private vertices gives
\[
 t(J^{[r]},U)=t(J,M_U),\qquad \mathbb E M_U=\mathbb E U=p.
\]
Even cycles are Sidorenko.  Both even and odd generalized theta graphs
are Sidorenko by Im--Li--Liu, Theorem~1.3 and Proposition~4.3,
respectively~\citep{imliliu2024}.  These graphon statements apply to a
finite weighted kernel by representing its states by intervals of the
corresponding weights.  Therefore $t(J,M_U)\ge p^m$, as required.
\end{proof}

This proposition concerns one atom.  It neither gives a rank-three
counterexample nor asserts closure under arbitrary one-point assembly.
In particular, the all-odd theta case cannot be discarded by imposing
an odd-order convention on the core.

\subsection{A finite sparsity certificate and sparse random hosts}

Let $J$ be the point--face incidence graph of OpenAI's $13$-point,
$22$-triangle complex~\citep[finite-complex proposition, part~(ii), and
exposure-order table]{openai161}.  The following balancedness calculation
does not use its rank-layer estimate.  We retain the face labels and
the two exposure orders of the source pinned in the bibliography:
\[
\begin{split}
 &(0,1,2,4,5,16,14,18,20,19,13),\\
 &(3,9,8,10,11,17,15,12,21,7,6).
\end{split}
\]
After the first triangle, the corresponding new points are, respectively,
\[
 (9,2,10,7,5,6,12,8,11,4),\qquad
 (4,3,11,10,5,6,8,12,1,7).
\]
These are directly checked membership and freshness certificates.
Every nonempty selection of $h$ faces from either order therefore uses
at least $h+2$ points: its first face uses three points, and every later
selected face retains its certified new point.

\begin{proposition}[Strict balancedness of the incidence example]
\label{density-incidence-strict-balance}
Every subgraph $Q\subseteq J$ with $v(Q)\ge3$ satisfies
$e(Q)\le2v(Q)-4$.  The graph $J$ is strictly $2$-balanced, with
\[
 m_2(J)=\frac{65}{33}.
\]
For every $r\ge2$, its private expansion $A_r=J^{[r]}$ is strictly
$r$-balanced, with
\[
 m_r(A_r)=\frac{65}{65r-97}.
\]
Here $m_r(H)=\max_{Q\subseteq H,\ v(Q)>r}(e(Q)-1)/(v(Q)-r)$.
\end{proposition}
\begin{proof}
Suppose $Q$ retains $a$ point nodes and $b$ face nodes.  If $a\ge2$,
at most $2a-4$ faces have all three points retained, by the two exposure
certificates.  Every other retained face contributes at most two
incidences, so $e(Q)\le2b+2a-4$.  When $a=1$, the condition $a+b\ge3$
gives $b\ge2$ and $e(Q)\le b\le2(a+b)-4$; when $a=0$, the assertion
is immediate.  Removing edges only decreases these bounds.

Since $v(J)=35$ and $e(J)=66$, its $2$-density is $65/33$.
If $3\le v(Q)<35$, then
\[
 \frac{e(Q)-1}{v(Q)-2}
 \le 2-\frac1{v(Q)-2}<2-\frac1{33}.
\]
A proper subgraph on all $35$ vertices has fewer than $66$ edges and
also gives a strict inequality.

Now select $f\ge2$ edges of $A_r$, and let $v$ be the number of
endpoints of the corresponding core edges.  Their union contains
exactly $v+(r-2)f$ vertices, since the private vertices are disjoint.
Any additional isolated selected vertices can only increase the
denominator.  The core estimate gives
\[
 \frac{f-1}{v+(r-2)f-r}
 =\frac1{r-2+(v-2)/(f-1)}
 \le\frac{65}{65r-97}.
\]
Strictness follows unless all core edges and vertices are retained
and no additional isolated vertex is selected.  A selection of one
edge with more than $r$ vertices has ratio zero; an edgeless selection
has negative ratio.  Thus neither affects the maximum.
\end{proof}

This finite calculation places the same atom family in the scope of
Nie--Spiro's random Tur\'an theorem~\citep{niespiro2025}.
Write $s(A_r)$ for their Sidorenko excess, namely the supremum of
$s\in\mathbb R$ for which some finite host $G$ has
$0<t(A_r,G)=p_r(G)^{66+s}$, with $0<p_r(G)<1$.
The supremum need not be attained.  Let $G^{(r)}_{n,\rho}$ denote the random
$r$-uniform hypergraph with edge probability $\rho$.
Set $\alpha=r-97/65$.  For fixed $r$ and
$\rho\ge n^{-\alpha}$, their theorem gives, asymptotically almost surely,
\[
 \operatorname{ex}(G^{(r)}_{n,\rho},A_r)
 \ge n^{97/65-o(1)}
       \bigl(\rho n^{r-97/65}\bigr)^{s(A_r)/(65+s(A_r))}.
\]
The balancedness is unconditional.  The conclusion $s(A_r)>0$,
and hence a positive gain factor when $\rho n^\alpha>1$,
uses the stated rank-layer and proper-subgraph-flat transfer input.
No explicit uniform lower bound on $s(A_r)$ or matching upper bound
is asserted here.
If $\rho n^\alpha\ge n^\varepsilon$ for a fixed $\varepsilon>0$,
positive excess gives a polynomial gain.  At the exact threshold the
displayed gain factor is one.

\subsection{Keeping all attachment ports}

Let the hyperedges of a finite uniform pattern $H$ be partitioned among pieces
$A_i$.  Let $P$ be its shared vertices and $I_i$ the vertices belonging
only to piece $i$.  Assume these sets give the full original vertex
decomposition $P\sqcup\bigsqcup_i I_i$; this is not inferred from an
edge partition alone.  Isolated vertices can be represented as a
free factor, or removed when probability weights are normalized.
The port set of $A_i$ consists of all its
vertices lying in $P$, not only one chosen root.  For a port assignment
$x$, put
\[
 \Phi_i(x)=\mathbb E_{y\in\Omega^{I_i}}
                 \prod_{e\in E(A_i)}U\bigl((x\sqcup y)|_e\bigr).
\]
The notation means that each factor is evaluated at the actual
endpoints of its edge.  Independent sampling uses the given finite
probability weights on $\Omega$.

\begin{proposition}[Exact multi-port contraction]
\label{density-multiport-contraction}
With these original edge and vertex decompositions,
\[
 t(H,U)=\mathbb E_{x\in\Omega^P}\prod_i\Phi_i(x|_{P\cap V(A_i)}).
\]
The identity includes empty products, zero-weight states and pieces
with several attachment vertices.  It does not require symmetry or
nonnegativity if the weights and kernel values are interpreted
algebraically rather than as probabilities and an ordering of the
input slots of every edge is fixed.
\end{proposition}
\begin{proof}
Fix the shared assignment.  The edge partition groups the factors
by piece, and the internal assignment carriers are pairwise disjoint.
The finite sum over their product therefore factors into the product
of the sums defining $\Phi_i$.  Averaging the shared assignment gives
the displayed formula.
\end{proof}

Thus unrooted atom densities are not sufficient attachment data.
The two-piece one-port case is the covariance identity already used
in the manuscript.  With several joints, a one-root summary can lose
correlations between ports of the same atom.  A natural remaining
question is whether an obligatory assembly whose individual atoms
are Sidorenko must itself be Sidorenko.  The contraction formula is
an exact starting point for this question, not a positive or negative
answer to it.

% The accompanying references.bib contains only the three directly cited
% primary sources. Im--Li--Liu is cited in its published CPC version;
% Theorem 1.3 and Proposition 4.3 retain the even/odd parity scopes used here.
% The Nie--Spiro formula was checked against arXiv:2309.12873, Theorem 1.4;
% the prose avoids transferring that preprint numbering to the publication.
